\begin{abstract} 
 Dark matter is thought to comprise the majority of the mass budget of the universe, with dark matter estimated to account for nearly $85\%$ of the total mass of the universe \cite{collaboration2020planck}. 
 Despite it's abundance, the nature of dark matter remains elusive.
 While Large scale cosmological probes such as Planck \cite{collaboration2020planck} and surveys of the spatial and mass function distributions of galaxies \cite{white1978core,white1991galaxy,de2008high,weinberg2015cold} have provided decisive evidence for an empirical model of dark matter that is cold and collisionless on the largest scale (CDM), the behavior of dark matter on the sub-galactic scale remains an area of active investigation and vigorous debate; one may hope that hints towards the underlying nature of dark matter can be found in this regime.
 A promising way to probe the physics of dark matter on the sub-galactic dark matter halos is by using anomalous flux ratios of quadruply lensed quasars. 
 In these systems, a light from a quasar passes through a central galaxy and is gravitationally lensed, producing four observable images. 
 Due to the non-linear nature of the dependence of flux ratios on the second derivative of the potential, information about dark matter on the sub galactic can be extracted from the flux ratios \cite{10.1093/mnras/stz3480}; several studies have used these methods to place constraints on Warm Dark Matter \cite{10.1093/mnras/stz3480, nadler2021dark} as well as self-interaction of dark matter \cite{2023PhRvD.107j3008G} without the need for luminous tracers required by many other dark matter probes.  
 In this work, study simulations and models of dark matter substructure of group mass ($\sim 10^{13} M_{\odot}$) to develop informed priors on dark matter models to support strong gravitational lensing studies.
 We split this in two sections: first we explore developing informed priors on the projected number density of subhalos within mass lensing halos, then we explore developing a new ultra fast simulation pipeline for simulating the evolution of dark matter halos in alternative dark matter models.
 
 Perturbations in the density of dark matter gravitationally collapse into objects known as dark matter halos \cite{press1974formation}; understanding how these perturbations grow and evolve is important for placing constrains on dark matter models.  
 In this work, we focus in particular on subhalos: halos evolving under the gravitational field of a host. 
 N-body simulations suggest that subhalos account for $5$--$20\%$ of the total number of halos near the lensed images in projection with the subhalo to isolated halo ratio depending strongly on redshift \citep{10.1093/mnras/stu2673, 10.1093/mnras/sty159}; although subhalos are subdominant in , \citet{gilman2024turbocharging} showed that uncertainty in tidal stripping of models is a major source of uncertainty in strong gravitational lensing studies. 
 Subhalos can be numerically challenging to study: subhalos can lose $>99\%$ of their mass due to tidal stripping, leaving the subhalo prone to artificial noise in N-body simulations and possibly artificially disrupting, which would fictitiously lower the number of subhalos predicted by these simulations {\color{red} CITE THE CDM REMMENANTS PAPER HERE}.
 To overcome these challenges we use the state-of-the-art semi-analytic \gal {\color{red} CITATION HERE} model tuned to ultra high resolution N-body simulations that minimize the effects of artificial disruption.
 Using this model we develop informed priors on the projected halo mass function expected by lensing studies, develop models for the scaling of the halo mass function with redshift and halo mass and explore the tidal effects of a realistically evolving potential of a central galaxy. 
 We find that a factor of $\sim 2$ accounts for theoretical uncertainties in the subhalos mass function and find that for current lesning studies the central galaxy has a currently negligible impact on the projected number density of substructure. 

 Recent inferences on the sub-galactic from direct gravitational imaging the potential of lenses \cite{vegetti2010detection, minor2021unexpected, tajalli2025sharp} have provided evidence for anomalously dense dark matter halos compared to CDM predictions.
 Although these results have been hotly debated and solutions within the CDM framework have been proposed \cite{he2025not}, the intriguing possibility remains that these results hint at physics beyond CDM; possible explanations include a class of alternate dark matter models with non-negligible self-interaction on the sub-galactic scales, called self-interacting dark matter (SIDM).
 In these models, dark matter self-interaction facilitates the conduction of heat within the halo; due to the negative heat capacity of self-gravitating systems, runaway heat transfer can occur leading to the halo gravothermally collapsing, leading to extremely dense collapsed remnants at the end of their evolution \cite{tulin2018dark,adhikari2025astrophysical}.
 3-D N-body simulations, which are traditionally the primary method of studying dark matter halos are computationally expensive especially in the core collapse regime; tidal mass loss amplifies this challenge by requiring a larger number of particles to resolve the halo initially, often requiring tens of thousands of core hours to properly resolve core collapse  even in idealized simulations. 
 In this work, we develop a pipeline to simulate the evolution of idealized SIDM subhalos in 1-D with realistic tidal evolution; by simulating subhalos evolution in 1-D, we are able to simulate SIDM physics much more computationally efficiently than the 3D case, resulting in over an order of magnitude speed up in some case enabling previously infeasible exploration of SIDM model space. 
 Using our models we explore the core collapse rates for a wide range for a wide range of SIDM cross sections for a realistic population of halos. 
\end{abstract}
